
For representation research the implication is narrower than a leaderboard result but harder to work around. In our lupus data pseudobulk beats frozen embeddings in strict source-only transfer, and added pooling capacity does not change that. The more durable point is that internal accuracy in these cohorts is not a measurement of the property that matters, and no amount of model development fixes a cohort in which design predicts the diagnosis. That is a recruitment and data-release problem. The fix is to record and distribute collection metadata, and to build cohorts in which cases and controls appear together within batches - not to build a better encoder.

None of these results supports a diagnostic claim. The external AUCs we report rank 261 or 26 donors and are not calibrated probabilities. What the analysis supports is a reporting discipline. A patient-level classifier should be published with three things. First, the design-only AUC of its development cohort \textbf{by variable group}, so a reader can see whether the association runs through collection, case mix or sample quality. Second, both permutation tests, or a plain statement that the collection-preserving one is undefined. Third, an external result on a cohort collected through a different route, fitted on the source cohort alone. None of that requires a new method. It requires releasing the metadata and running two cheap checks.

\section{8. Limitations}\label{limitations}

These checks see only recorded design. Unmeasured collection structure - referral patterns, institutional treatment protocols, sample handling not captured in released metadata - is invisible to every quantity we report. A low design-only AUC is therefore weak evidence of an unaffected cohort rather than strong evidence of one. The Ren cohort shows the same failure in its recorded form: our first set of design variables missed the recruiting hospital entirely and reported a collection-group AUC of 0.553 where the completed set gives 0.750.

\textbf{The collection-preserving permutation conditions on a subset of the recorded metadata.} Its strata come from collection variables; sample-quality and demographic columns are in the design matrix but not in the strata. Section 3.4 measures the consequence directly: with a sample-quality variable still associated with the diagnosis inside a stratum and driving expression, the test rejects far above nominal with no disease effect anywhere in the generating model. A significant collection-preserving result is therefore evidence against the collection strata, not evidence of a disease-specific signal, and it is not a demonstration that a ``design-independent component'' exists. Building a null that conditions on the full matrix, including its continuous columns, would require a conditional randomisation scheme we do not develop here. That is the single largest methodological gap in this paper.

\textbf{The three variable groups are not equivalent, and only one is study design in the strict sense.} Demographics can be case mix or genuine risk; sample-quality summaries sit downstream of the assay but can also sit downstream of the disease. We therefore report the collection block separately everywhere and treat it as primary for any claim worded as being about study design. The headline count differs between the two definitions, and both are given.

V\_D is linear by construction and depends on how the design matrix is coded. We reduce but do not remove this by reporting nonlinear design-only classifiers, the variable grouping, and the design matrix's own null. The nonlinear arms use two tree ensembles and are not an exhaustive comparison of adjustment algorithms.

\textbf{We cannot fully account for the CMV residualisation loss.} Section 3.5 shows that design-matrix width and raggedness produce a real loss with no confounding present, but less than the 0.282 observed. Something about that cohort's design matrix that we have not isolated accounts for the remainder.

All nine comparisons were run at five seeds, and every value we quote is the range over those five. Three comparisons have fewer than 15 donors in the minority class. We flag rather than exclude them, but no conclusion should rest on them.

The nine comparisons come from five datasets and are \textbf{not independent}. Comparisons drawn from the same dataset share donors, so the count of eight in nine is a count of comparisons, not of independent replications.

Matched random subsets separate sample-size loss from restriction without identifying a causal technical effect, so the restriction results show non-attributability rather than that retained signal is biological. Our coverage is four diseases from five cohorts, all peripheral blood, all public, all curated under one schema. Tissue-resident and non-blood settings are untested. Geneformer pretraining overlap with individual benchmark cells was not tested.

Finally, these are checks. They order cohorts and point at variable groups. They do not identify a causal effect, and we make no claim that they do.

\section{Methods}\label{methods}

\section{9. Materials and methods}\label{materials-and-methods}

\subsection{9.1. Cohorts and the unit of analysis}\label{cohorts-and-the-unit-of-analysis}

Seven public datasets were used, in two groups. Table S6 is the registry. It gives, for each dataset, the identifier, the donor count reported by the source, and the donor, case and control counts actually modelled - these differ, because donors below the minimum cell count or carrying a label outside the contrast are dropped. Five entries were verified against the CELLxGENE Discover API on 2026-09-07; the two used only in Section 5 come from GEO and are marked as not API-verified.

\textbf{Five datasets enter the main screen}, all obtained through the CELLxGENE Discover curation API: the CD4-positive alpha-beta T-cell subset of GSE174188 (261 donors; 162 cases, 99 controls) {[}2{]}, the Ren COVID-19 atlas (196 donors) {[}39{]}, the Stephenson COVID-19 cohort (120 donors) {[}40{]}, the COMBAT consortium (124 donors) {[}41{]} and the HIHA cytomegalovirus cohort (108 donors) {[}42{]}. Nine case-control comparisons are formed from these five, so the comparisons are not independent of one another; Table 1 gives the donor overlap.

\textbf{Two further lupus datasets are used only in the downstream analyses of Section 5}, not in the screen: GSE135779 (44 childhood donors; 33 cases, 11 controls) {[}1{]} and GSE285773 CD4-positive T cells (26 donors; 16 cases, 10 controls) {[}3{]}, the latter as the independent transfer target.

Registry entries, dataset identifiers, donor counts and download sizes were verified against the API on 2026-09-07 and are released in \texttt{cohorts/registry.yaml}.

The donor was the independent unit in every supervised analysis. Cells from one donor never crossed a training/test partition.

\subsection{9.2. Donor-level tables and the design variables}\label{donor-level-tables-and-the-design-variables}

Each cohort was streamed from its h5ad in 100,000-cell chunks and reduced to one row per donor. Expression was aggregated to log1p CPM pseudobulk. The layer used for aggregation was chosen by sniffing for integer counts rather than assuming a layer name.

Design variables were drawn from schema-guaranteed observation fields (\texttt{donor\_id}, \texttt{disease}, \texttt{sex}, \texttt{development\_stage}, \texttt{assay}, \texttt{tissue}) plus cohort-specific collection variables detected by pattern from the distributed object. Detection patterns cover batch, pool, run, lane, chip, site, centre, city, region, province, hospital, clinic, study, sub-study, source, dataset and donor source. Ages given as strings (``25-year-old stage'', ``fifth decade stage'', ``90 year-old and over stage'') were parsed to years.

Covariates that are consequences of the disease were excluded by an explicit blocklist matching severity, outcome, comorbidity, medication, sample timing, symptom, diagnosis, stage, grade and WHO score. Supplementary Table S2 lists, for every cohort, which columns entered the design matrix, in which group, and which were excluded and why.

Variables were assigned to three groups. \textbf{Sample quality}: cells per donor, mean UMI per cell, mean genes per cell, aggregate mitochondrial percentage, and their logs. \textbf{Demographic}: age in years, sex, ethnicity. \textbf{Collection}: every detected \texttt{batch\_\_*} variable, assay and suspension type.

\subsection{9.3. Residual label variance}\label{residual-label-variance-1}

For centred diagnosis vector \emph{y} and recorded design matrix \textbf{D} including an intercept, V\_D = 1 − R²(y \textasciitilde{} D). Two versions are reported. The in-sample version uses ordinary least squares on all donors and is reported for comparison. The cross-fitted version replaces the fitted values with out-of-fold predictions from five-fold \texttt{KFold} with shuffling, seeded by the run seed, and is the version used everywhere a number is interpreted. Negative R² values were clipped at zero before subtraction.

\subsection{9.4. The design matrix's own permutation null}\label{the-design-matrixs-own-permutation-null}

For each design matrix, the diagnosis vector was permuted 200 times. Each permuted label was pushed through the \textbf{same five folds} used for the observed value, so the null is a distribution of the cross-fitted statistic, not of a different one. The reported p-value is one-sided, (1 + \#\{null ≤ observed\}) / (200 + 1), and answers: how often does a random diagnosis leave at most as much unexplained variation as the observed one? The null mean and its 2.5th percentile are reported alongside.

The design-only AUC is tested the same way, with 200 permutations of the diagnosis through the identical frozen pipeline (\texttt{p\_design\_auc} in Table S4); that is the primary test, for the reason given in Section 2.2. The in-sample V\_D is released with its own in-sample null, as \texttt{V\_D\_insample} and \texttt{p\_V\_D\_insample} in Table S4. The two are separate quantities and are never mixed: p(V\_D) always refers to the cross-fitted pair. This null is what makes either version readable, because its location depends on the number of design columns relative to donors.

\subsection{9.5. Design-only and diagnosis classifiers}\label{design-only-and-diagnosis-classifiers}

Numeric design variables were median-imputed and standardised inside each outer training fold; categorical variables were most-frequent-imputed and one-hot encoded, with unseen held-out levels ignored. Every grid and fixed setting is listed in Table S7. All imputation, encoding, scaling, feature selection and regularisation choices were fitted on training donors only.

The main classifier was balanced logistic regression (liblinear, \texttt{max\_iter} 20,000), with inverse regularisation strength selected from 10⁻⁴ to 10⁴ in decade steps by five-fold resampling inside each outer training set. Internal analyses used 20 repeated stratified five-fold donor splits with integer seeds 20260801-20260820. Out-of-fold predictions were pooled within each repeat; we report the mean AUC and the 2.5th-97.5th percentile of repeat-level AUCs. That range is a descriptive split-sensitivity range. It is not a confidence interval, and it does not treat correlated folds or repeats as independent {[}31{]}.

\textbf{Nonlinear design-only arms.} We used two tree ensembles rather than a wider search because the question is whether a \emph{low-dimensional} design-diagnosis relationship is being missed by a linear model, not which learner maximises AUC. Random forests used 256 trees, maximum depth 4, minimum leaf size 3. Histogram gradient boosting used 150 iterations, learning rate 0.05, seven terminal leaves, minimum leaf size 5. Both used fixed hyperparameters, fitted inside every outer training fold, and were not selected against held-out donors. Reporting them with fixed settings is deliberate: a tuned nonlinear arm would not be comparable with the tuned linear one under the same permutation null.

\subsection{9.6. The two permutation tests}\label{the-two-permutation-tests}

Both permutation tests and the observed statistic they are compared against run one identical frozen pipeline: standardisation, PCA to 50 components (or fewer when donors are limiting), balanced logistic regression at fixed inverse regularisation 1.0, and a single cross-fitting repeat. Freezing is necessary: comparing a tuned observed statistic against an untuned null biases p-values downward.

The \textbf{free} test permutes the diagnosis vector without restriction. The \textbf{collection-preserving} test permutes it only within collection strata. A stratum is the interaction of every detected \texttt{batch\_\_*} variable with assay, where assay varies. Table 1 names the stratum variables for every comparison, and the screen writes them into its own output (\texttt{strata\_definition}) so the name and the implementation cannot drift apart.

\textbf{The conditioning set is collection variables only.} Sample-quality and demographic columns are in the design matrix but not in the strata, so this test conditions on a subset of the recorded metadata. Section 3.4 measures the consequence: with a sample-quality variable still associated with the diagnosis inside a stratum and driving expression, the test rejects at a rate well above nominal even though the generating model contains no disease effect. Conditioning on the full matrix would need a conditional randomisation scheme for continuous columns, which we do not attempt here.

Where a cohort records no collection variable at all, the fallback is tertiles of log cells per donor. That is a sample-quality variable, so for those comparisons the test is a sample-quality-preserving permutation and is labelled as such rather than being folded in silently.

A permutation contributes to the collection-preserving null only if at least one stratum contains more than one donor and both labels; if no permutation qualifies, the p-value is undefined and reported as such rather than as a number. Main runs used 1,000 permutations, so the smallest value it can return is 1/1001 ≈ 0.0010. The V\_D and design-only AUC nulls use 200, so their floor is 1/201 ≈ 0.0050. Every seed used the same counts; the two floors in Table 1 are these two tests, not two different runs.

\subsection{9.7. Feasibility and power checks}\label{feasibility-and-power-checks}

A comparison is reported as \textbf{not answerable with these data} if the collection-preserving permutation has no qualifying stratum, or if the design-only AUC reaches 1.000, which is the complete-collinearity case. It is reported as \textbf{underpowered}, and kept in the main ordering with a flag, if it has fewer than 40 donors or fewer than 15 in the minority class. The two are separated because they call for different responses: the first cannot be fixed by more careful analysis, and the second can be fixed by more donors.

The difference between the tuned and the frozen pipeline AUC is reported as a column (\texttt{frozen\_minus\_tuned\_auc}, Table S4) and is deliberately \textbf{not} used as a feasibility rule. A gap between two different algorithms does not establish that a comparison is unanswerable. Every inferential statement in this paper is made about the frozen pipeline, which is the statistic both permutation nulls are built from; the tuned AUC is descriptive.

\subsection{9.8. Seed stability}\label{seed-stability}

All nine comparisons were re-run at five seeds (20260907-20260911), varying the seed for fold construction, permutation draws, and the cross-fitting split used for V\_D and its own permutation null. All five per-seed outputs are released rather than summarised (Figure S7, Table S8), and every range quoted in Table 1 and in the text is the minimum and maximum over those five runs. During this work we found a bug: the seed was bound as a default argument at function definition time, so \texttt{-\/-seed} never reached the V\_D null and the V\_D columns came out identical across seeds. The corrected implementation reads the seed at call time. Every number reported here comes from the corrected version.

\subsection{9.9. Simulation}\label{simulation}

Cohorts of 200 donors and 100 features were generated. A latent variable \emph{z} produced a site indicator and a continuous quality proxy; the diagnosis was generated as ρ·z + √(1−ρ²)·e and then either used continuously or dichotomised. Biological and technical loading vectors overlapped partially (technical = 0.7·random + 0.3·biological, renormalised). Seven regimes were crossed with ρ ∈ \{0, 0.25, 0.5, 0.75, 0.9, 0.98, 1\} and both label types, at 200 replicates per cell, giving 19,600 simulated cohorts. The seven regimes are:

\begin{enumerate}
\def\labelenumi{\arabic{enumi}.}
\tightlist
\item
  \textbf{additive linear}, the base case;
